\section{Introduction}
\label{sec:introduction}

Drought frequency and severity have increased worldwide, driven by
ongoing climate change and rising anthropogenic water demands
\citep{IPCC2023, Pokhrel2021}, placing growing pressure on surface water
storage systems. In the Mediterranean basin, these trends are particularly
acute: multi-year dry spells have repeatedly pushed reservoir stocks to
critically low levels, threatening agricultural supply, hydropower
generation, and urban water distribution \citep{Cammalleri2022,
Bevacqua2023}. Sicily offers one of the most illustrative expressions of
this vulnerability: over the past decade, its reservoirs have repeatedly
reached critically low levels during the dry season \citep{ISPRA2023},
and total stored volumes have rarely recovered to historical averages even
after rainfall events, signalling a progressive depletion trend in
regional water reserves.

Reservoir storage monitoring depends on two observables, water level and
surface area, and on how reliably and often each can be tracked. Satellite
radar altimetry, from traditional nadir missions to the newer wide-swath
SWOT mission \citep{Biancamaria2016}, feeds databases such as DAHITI
\citep{Schwatke2015} and GROWL \citep{Zhang2026} that extend level monitoring
beyond the world's sparse in-situ gauge network, yet even SWOT revisits any
one reservoir only every $\sim$21 days. Surface area, by contrast, can be
measured by any imaging sensor whose footprint overlaps the reservoir, far
more often.

Converting either observable into a storage estimate requires an
area--elevation--volume (AEV) relationship, and databases such as GRanD
\citep{Lehner2011} and GRDL \citep{Hao2024} already compile one for many
reservoirs, built from local bathymetric or design surveys. For a
significant number of reservoirs, however, these AEV relationships need to
be updated or recreated from scratch. Building or updating them by
combining altimetry and area observations, as proposed by
\citet{Schwatke2020}, is one route toward closing that gap. Its payoff
would be substantial: a satellite-derived AEV relationship would extend
storage monitoring to the many reservoirs no recent survey has reached,
but only on top of an area (or level) estimate that is itself automated
and reliable at scale. Establishing that reliability for SAR-derived area
is our focus here.

Among the sensors capable of that area retrieval, SAR is particularly well
suited: it delineates the water surface extent directly, and unlike optical
sensors operating
in the visible and near-infrared range, SAR is not affected by cloud cover
or solar illumination \citep{Torres2012, Mullissa2021}, a critical
advantage in Mediterranean and tropical contexts where cloud contamination
can severely reduce the temporal availability of optical observations
\citep{Pekel2016}. Processed on cloud computing platforms such as Google
Earth Engine (GEE), Sentinel-1 SAR data enable systematic surface water
mapping over large areas and long time periods \citep{Gorelick2017}.
Recent studies confirm this operational maturity. \citet{BhatpuriaEtAl2024}
applied automated SAR workflows for storage monitoring in semi-arid India.
\citet{ValerioEtAl2024} showed that Sentinel-1 combined with Sentinel-2 on
GEE reliably detects water bodies smaller than 0.5~ha, and
\citet{PengEtAl2025} demonstrated GEE-based SAR--optical fusion for flood
mapping with accuracies above 92\%. Beyond this all-weather advantage,
area-based SAR masks are spatially explicit and can capture asymmetric
drawdown patterns and disconnected water pools that a single water-level
number cannot.

Despite this progress, SAR-based reservoir delineation exhibits
substantial variability in performance across water bodies that cannot be
attributed to sensor characteristics alone. Wind-driven surface roughness
raises backscatter over open water, causing misclassification as non-water
\citep{Bioresita2018, PenaLuque2021}. At the water--land boundary,
mixed pixels introduce systematic errors whose severity depends on the
ratio of shoreline length to water body area \citep{PenaLuque2021}.
Shadowing and layover effects in topographically complex settings further
complicate delineation near steep reservoir margins \citep{Martinis2015}.
These phenomena are well-established individually, yet a basic operational
question remains unanswered: \emph{which geometric property of a reservoir best
predicts its SAR classification skill, knowable in advance of any SAR
processing, and does that skill in fact require the added complexity of a
dual-polarisation, machine-learning classifier, or does a simple
single-polarisation threshold already suffice?}

We propose that the answer lies in the \textbf{area-to-perimeter ratio
(A/P, in metres)}, a length-scale index that captures the compactness of
the shoreline relative to the water body's spatial extent. The physical
reasoning proceeds as follows. For a sensor with ground resolution $r$,
the boundary strip of width $r$ enclosing the perimeter $P$ contains
approximately $P \cdot r$ square metres of ambiguous pixels (neither
purely water nor purely land). As a fraction of total water area $A$,
this amounts to $P \cdot r / A = r / (A/P)$, a quantity that grows as
$A/P$ decreases, that is, as the shoreline becomes more irregular or
elongated relative to the enclosed area. This geometric argument is
corroborated by \citet{PenaLuque2021}, who quantified a systematic
perimeter-inward erosion of $\sim$20~m in Sentinel-1-derived water masks
at the pixel level. Because this mixed-pixel argument depends only on
shoreline length and pixel size, it is sensor- and algorithm-agnostic, and
predicts that A/P is a reliable a-priori indicator of
the \emph{geometric} ceiling on SAR classification performance, one that,
to our knowledge, has not been systematically tested across a multi-reservoir
dataset. Radiometric failures form a second, orthogonal axis to which A/P is
blind: wind-driven roughening in particular raises open-water backscatter and
can trigger omission errors, and has been proposed to scale with wind fetch
and hence with reservoir elongation \citep{McQuillan2025}. Whether this
radiometric axis genuinely co-varies with geometry or acts independently is an
empirical question we test directly rather than assume.

This paper addresses that gap with a fully automated, cloud-based GEE
workflow applied to 62 reference-quality-screened reservoirs on five
continents, referenced against the
JRC Global Surface Water dataset \citep{Pekel2016}, and to four Sicilian
reservoirs validated against high-resolution PlanetScope imagery as near-truth.
First, we quantify the statistical relationship between A/P and SAR
surface-area accuracy across the global set, the geometric half of the
question above. Second, we compare a single-polarisation (VV) per-scene
threshold against a dual-polarisation (VV+VH) per-scene SVM under an
identical pipeline, testing whether the added complexity of the dual-band
classifier is justified on accuracy grounds. Added complexity is usually
assumed to come at an added computational cost too, so we audit that cost for
each detector as well. Because training samples are drawn automatically from
global water masks, the workflow carries no manual delineation and transfers
directly across regions. By testing two classification methods across many
reservoirs, our central objective is to obtain a parameter, computable from
global databases such as GRanD or GROWL before any SAR image is processed,
that indicates in advance how reliable the resulting water mask will be.

The remainder of the paper is organised as follows.
Section~\ref{sec:study_area} describes the two-tier reservoir sample.
Section~\ref{sec:methods} details the SAR preprocessing, the automated
training, the two scene-adaptive classifiers, the morphometric analysis, and
the validation and cost-audit strategy. Section~\ref{sec:results} presents the
A/P--performance relationship, the single- versus dual-polarisation
comparison, the computational-cost audit, and the secondary climate and
reference limits.
Section~\ref{sec:discussion} discusses implications for operational monitoring,
the geometric-versus-radiometric limits of the approach, and avenues for
future work, including the conversion of surface area to storage volume.
Section~\ref{sec:conclusions} summarises the main conclusions.
